\documentclass[11pt,a4paper]{article}

\usepackage{geometry,microtype}
\usepackage{amsmath,amssymb}
\usepackage{booktabs}
\usepackage{enumitem}
\setlist{nosep}

\usepackage[shore]{../lrssntheme}
\usepackage{hyperref}
\hypersetup{colorlinks=true,linkcolor=LrssnPrimaryText,
            citecolor=LrssnPrimaryText,urlcolor=LrssnPrimaryText}

\geometry{left=25mm,right=25mm,top=24mm,bottom=24mm}
\setlength{\parindent}{0pt}
\setlength{\parskip}{0.6\baselineskip}

\makeatletter
\renewcommand{\maketitle}{%
  \begin{center}
    {\displayfont\LARGE\bfseries\color{LrssnPrimaryText}\@title\par}
    \vspace{0.5em}
    {\large\@author\par}
    {\small\color{LrssnMuted}\@date\par}
  \end{center}
  \vspace{1em}
}
\makeatother

\title{Short Paper Title}
\author{Author Name}
\date{\today}

\begin{document}
\maketitle

\begin{abstract}
\noindent
Write a brief summary of the purpose, method, and main result.
\end{abstract}

\section{Introduction}
Introduce the topic and state the question addressed by the paper.

\section{Main content}
Present the central argument, observation, or method.

\subsection{Example point}
Use subsections to organise a longer short paper.

\begin{equation}
    y = mx + c
\end{equation}

\section{Conclusion}
Summarise the conclusion and, if useful, identify a next step.

\end{document}
